\documentclass[11pt,letterpaper]{article}

% ======= PACKAGES =======
\usepackage[utf8]{inputenc}
\usepackage{geometry}
\geometry{margin=1in}
\usepackage{setspace}
\setstretch{1.5}
\usepackage{fancyhdr}
\usepackage{xcolor}
\usepackage{titlesec}
\usepackage{lastpage}
\usepackage{hyperref}
\usepackage{mathpazo}
\usepackage{graphicx}
\usepackage{subcaption}
\linespread{1.5}

% ======= HEADER & FOOTER =======
\pagestyle{fancy}
\fancyhf{}

% Header (faded + smaller)
\fancyhead[L]{\footnotesize\textcolor{gray!90}{\textit{Project Description}}}
\fancyhead[R]{\footnotesize\textcolor{gray!90}{\textit{Mostafa Rezaali}}}

% Footer (faded)
\fancyfoot[L]{\footnotesize\textcolor{gray!90}{\textit{NSF AGS-PRF -- University of Florida}}}
\fancyfoot[R]{\footnotesize\textcolor{gray!90}{Page \thepage}}

\renewcommand{\headrulewidth}{0.4pt}
\renewcommand{\footrulewidth}{0.4pt}

% ======= PARAGRAPH STYLE =======
\setlength{\parindent}{15pt}
\setlength{\parskip}{6pt}
\titlespacing*{\section}{0pt}{0.45em}{0.15em}
\titlespacing*{\subsection}{0pt}{0.35em}{0.1em}

% ======= DOCUMENT BEGINS =======
\begin{document}

% ======= TITLE BLOCK =======
\begin{center}
    \vspace*{-1.5cm}
    {\Large\bfseries Postdoctoral Fellowship Application} \\[0.2cm]
    {\small\textit{NSF Atmospheric and Geospace Sciences Postdoctoral Research Fellowship}} \\[0.05cm]
    {\small\textit{Department of Geography, University of Florida}} \\[0.1cm]
    \rule{0.6\textwidth}{0.4pt}
\end{center}
\vspace{-0.5cm}

\noindent
Mostafa Rezaali \\
Ph.D. Candidate, Climate Science (Geography) \\
University of Florida \\
mostafarezaali@gmail.com \quad $\mid$ \quad (352) 709-5350 \\[0.3cm]
{\large\bfseries Project Description}\\[-0.15cm]
\rule{\textwidth}{0.4pt}

\section*{Project Title}
\textbf{Postdoctoral Fellowship: AGS-PRF: Probabilistic Prediction of Heat Waves and Flash Droughts Using Physics-Informed Conditional Flow Matching on an Icosahedral Mesh}

\section*{Motivation and Central Hypothesis}
Heat waves and flash droughts are compound hydroclimate hazards with large consequences for public health, energy demand, water resources, ecosystems, and agriculture. Their impacts often emerge on decision horizons of two to four weeks, when emergency managers, utilities, farmers, and public-health agencies still have time to act, but when deterministic forecast skill is often weak. Predictability at these leads is not purely local. It reflects slowly evolving ocean-atmosphere patterns, stationary waves, soil-moisture memory, radiative seasonality, synoptic persistence, and land-surface feedbacks.

My central hypothesis is that a graph-based conditional flow model, trained directly on high-resolution U.S. temperature anomalies while conditioned on global circulation and teleconnection states, can reveal and exploit a measurable but spatially heterogeneous week-3 signal for extreme heat and flash-drought precursors. The proposed work treats AI not as a replacement for physical reasoning, but as a tool for testing where atmosphere-land-ocean memory produces useful probabilistic predictability.

\section*{Preliminary Work}
I have implemented MeshFlowNet, a GraphCast-style icosahedral mesh encoder-processor-decoder for direct 15-day CONUS daily maximum temperature prediction. The model uses current and lagged local PRISM temperature anomalies, atmospheric and land-surface fields, topography, latitude, longitude, seasonal insolation, land mask, five teleconnection indices, and a 59-channel global ERA5/coarse atmospheric context stack. It supports deterministic residual prediction and probabilistic conditional-flow-matching modes.

The current deterministic 5-fold leave-years-out hindcast covers MJJAS 1981--2023, with each year predicted by a model that did not train on that year. Across 5,332 samples, stitched model temporal anomaly correlation (TAC) is 0.1016 compared with 0.0735 for persistence, a +0.0281 improvement. Per-fold validation $R^2$ is approximately 0.58--0.60, and sample maps show skillful reconstruction of several high-amplitude heat anomalies. These results are preliminary, but they show that the model extracts nontrivial anomaly signal at a difficult lead time and justifies a systematic postdoctoral investigation of predictability, uncertainty, and physical interpretation.

\section*{Aim 1: Build a Calibrated Conditional-Flow Forecast System}
The first aim will convert the current deterministic residual model into a probabilistic forecast system. The deterministic branch predicts $y(t+15)$ as persistence plus a learned residual in one forward pass. The conditional-flow branch will learn a velocity field from initialization-time state to the verifying anomaly field, conditioned on teleconnection indices and global atmosphere/ocean fields. Ensembles will be generated by sampling the flow and by training multiple seeds per cross-validation fold.

Calibration will be assessed with continuous ranked probability score (CRPS), reliability diagrams for exceedance of local 90th and 95th percentiles, rank histograms, Brier score, and spatially stratified spread-skill relationships. This aim will test whether conditional flow matching can deliver useful uncertainty rather than only sharper deterministic maps. A successful result will be a reproducible probabilistic benchmark for week-3 U.S. heat prediction.

\section*{Aim 2: Diagnose Sources of Subseasonal Predictability}
The second aim will use controlled ablations and interpretable diagnostics to determine which features carry week-3 skill. I will compare models trained with: local persistence only; local meteorology plus static and seasonal terms; teleconnection indices; global ERA5 fields; and the full conditioning stack. The analysis will produce per-pixel TAC, skill over persistence, regional skill summaries for NOAA climate regions, block-bootstrap significance tests by year, and event composites for strong teleconnection regimes.

The scientific target is not only a higher metric, but a physical account of where AI forecast skill is consistent with known pathways such as Pacific-North American wave trains, moisture transport, soil-moisture feedback, subtropical ridge variability, and persistent high-pressure anomalies. These diagnostics will help distinguish robust predictive signal from overfitting or spatially smooth climatological reconstruction.

\section*{Aim 3: Connect Heat Prediction to Flash-Drought-Relevant Risk}
The third aim will connect temperature-anomaly prediction to flash-drought precursors by evaluating rapid intensification of heat and atmospheric demand in conjunction with antecedent soil-moisture state. The model already ingests soil moisture, humidity, winds, sea-level pressure, radiation, and large-scale circulation fields. I will evaluate whether the learned representations anticipate co-occurring hot, dry, high-demand regimes over agriculturally important regions.

Products will be reported as calibrated probabilistic exceedance maps, regional time series, and case studies rather than deterministic binary declarations. This framing keeps the science focused on predictability and uncertainty while making the outcomes relevant to water, agriculture, energy, and health decision contexts.

\section*{Data and Model Architecture}
The hindcast system will use PRISM daily maximum temperature over CONUS at 0.04-degree resolution, ERA5 atmospheric variables and global context fields, teleconnection indices, topography, latitude, longitude, land mask, top-of-atmosphere insolation, and seasonal encodings. The current local input stack contains current and lagged PRISM temperature fields; geopotential, soil moisture, sea-level pressure, 2-m temperature, 850-hPa humidity, 850-hPa temperature and winds, 300-hPa geopotential; topography; latitude and longitude; day-of-year sine and cosine; insolation; and land mask. The global context stack contains SST, OLR, winds, geopotential, temperature, humidity, surface pressure, sea-level pressure, and total column water vapor across multiple levels.

MeshFlowNet maps the CONUS grid to an icosahedral mesh, performs message passing with graph interaction networks and conditioning, and decodes back to the land grid. This representation preserves global-to-regional context while avoiding purely convolutional locality. The probabilistic branch will use continuous-time embeddings and conditional flow matching; the deterministic branch will remain as a transparent baseline and an efficient ablation.

\section*{Evaluation and Validation}
All claims will be based on leave-years-out cross-validation over 1981--2023. Metrics will include TAC, spatial anomaly $R^2$, MAE, MSE skill relative to persistence, CRPS, reliability of extreme exceedance probabilities, and block-bootstrap confidence intervals. Baselines will include climatology, persistence, ridge regression on teleconnection/global summary features, and published operational S2S reference skill where metric and domain are comparable. I will report negative or regionally limited results explicitly because understanding where AI models fail is essential for trustworthy climate-risk use.

\section*{Intellectual Merit}
The intellectual merit is the integration of subseasonal predictability theory, graph neural weather models, and probabilistic flow matching. The project will advance knowledge by quantifying when global and land-surface information improves 15-day U.S. heat prediction beyond persistence; producing spatially explicit diagnostics of learned teleconnection skill; and testing whether conditional-flow models can provide calibrated uncertainty for regional extreme-temperature risk. The outcome will be a physically evaluated AI framework rather than a black-box score table.

\section*{Broader Impacts}
The broader impacts center on usable, transparent climate-risk intelligence and training. First, the project will produce reproducible code, derived hindcast diagnostics, and documentation suitable for students and practitioners who need to understand AI forecast limitations. Second, I will develop a short teaching module on graph neural weather prediction and uncertainty for graduate climate modeling or environmental data-science courses. Third, I will mentor undergraduate and REU students in open climate-data workflows, with attention to students entering atmospheric science from geography, environmental engineering, and data science. Fourth, the forecast diagnostics will be framed around heat-health, water, energy, and agriculture use cases where calibrated uncertainty can be more valuable than a single deterministic forecast.

\section*{Host Institution, Mentoring, and Career Development}
The proposed host context is the Department of Geography at the University of Florida, which provides a strong intellectual environment for climate science, hydroclimatology, climate extremes, geospatial analysis, and AI-enabled environmental prediction. The fellowship will support a transition from doctoral model development to an independent research program in probabilistic subseasonal predictability, uncertainty quantification, and open forecast diagnostics.

Professional development will include regular research meetings, manuscript and proposal-development milestones, conference presentations, student-mentoring opportunities, training in responsible and human-centered GeoAI, and explicit preparation for independent faculty or research-scientist positions. The mentoring plan will emphasize scientific rigor, technical independence, communication to interdisciplinary audiences, and development of a visible research identity in AI-enabled atmospheric predictability.

\section*{Facilities and Resources}
The project will use research computing resources, Python/PyTorch software environments, geospatial and climate-data analysis tools, and publicly available PRISM/ERA5-derived data workflows. The local codebase already runs on GPU-equipped computing systems and includes scripts for training, hindcast export, bootstrap statistics, baseline comparison, TAC maps, sample maps, and CRPS/reliability analysis. The separate Facilities, Equipment and Other Resources upload states: See the Project Description, as requested by the AGS-PRF solicitation.

\section*{Timeline and Milestones}
Months 1--6: finalize data audit, deterministic baseline, ridge/climatology/persistence baselines, and per-pixel skill maps. Months 7--12: train probabilistic conditional-flow ensembles, compute CRPS and reliability diagnostics, and prepare the first manuscript on week-3 heat predictability. Months 13--18: complete teleconnection-regime ablations, block-bootstrap significance tests, and flash-drought-relevant risk diagnostics. Months 19--24: release reproducible workflows and derived hindcast products, submit a second manuscript, present results at AGU/AMS or comparable meetings, and complete career-development milestones with mentors.

\section*{Assessment of Success}
Success will be assessed by scientific and broader-impact outputs: (1) a documented cross-validated hindcast benchmark with uncertainty metrics; (2) at least two manuscripts, one on model and predictability and one on physical-regime interpretation or flash-drought-relevant risk; (3) public release of code, metadata, and nonrestricted derived products; (4) student mentoring and teaching materials; and (5) a clear professional transition from doctoral model development to an independent research program in AI-enabled atmospheric predictability.
\end{document}
